\section{The Sacraments at the Service of Communion}
\label{sec:communion-introduction}

Grace never isolates the believer. Though every soul stands singly before God in Baptism, Confirmation, the Eucharist, Penance, and the Anointing of the Sick, Christ did not die to gather scattered solitaries; He shed His Blood to form a living Body, a royal priesthood, and a household of faith (1~Pet~2:9; Eph~2:19). St.~Thomas Aquinas grounds the complete number of the seven sacraments in this twofold truth about human nature and supernatural life: man is perfected not only in his own person, through birth, growth, nourishment, and healing, but also as a member of a living society whose continuity requires both sacred governance and fruitful generation.
\begin{quotation}
In regard to the whole community, man is perfected in two ways. First, by receiving power to rule the community and to exercise public acts: and corresponding to this in the spiritual life there is the sacrament of order, according to the saying of Heb.~7:27, that priests offer sacrifices not for themselves only, but also for the people. Secondly in regard to natural propagation. This is accomplished by Matrimony both in the corporal and in the spiritual life: since it is not only a sacrament but also a function of nature. (\emph{ST}~III, q.~65, a.~1)
\end{quotation}
Holy Orders and Matrimony stand together at the completion of the sevenfold sacramental organism as the \textbf{Sacraments at the Service of Communion} (CCC~1533--1535). Neither is required of every individual Christian for personal salvation; both are indispensable to the enduring life and mission of the Church across history.

\subsection{Consecration for the Salvation of Others}
\label{subsec:communion-consecration}

Baptism, Confirmation, and the Eucharist consecrate the believer into Christ's own life; Penance and Anointing restore that life when wounded by post-baptismal sin or pressed by grave bodily infirmity. Holy Orders and Matrimony turn an already-baptized Christian outward toward a particular mission for the salvation of others. They confer genuine sanctifying grace upon the ordained minister and upon Christian spouses, yet they sanctify the recipient precisely through a lifelong, self-giving office in the Body of Christ (CCC~1534).

Viewed through the hierarchical theology of Pseudo-Dionysius, which St.~Thomas adopts in his reply on the sevenfold number (\emph{ST}~III, q.~65, a.~1, ad~3), every sacred economy comprises ministers who act, rites performed, and recipients who draw near. Holy Orders supplies the consecrated agents of the Church's liturgy, teaching, and pastoral rule; Matrimony brings forth the new members who are carried to the font and formed in faith to approach the altar. Viewed medicinally, as remedies against the fractures introduced by Adam's fall, Holy Orders defends the community against schism, division, and doctrinal ignorance, while Matrimony heals disordered desire and repairs the diminishment wrought by death (\emph{ST}~III, q.~65, a.~1).

In the Church's concrete life, these two sacraments build up the same Mystical Body along two complementary planes taught by the Second Vatican Council (\emph{Lumen Gentium}~10--11). Through Holy Orders, Christ perpetuates the apostolic hierarchy and the ministerial priesthood (\latin{sacerdotium ministeriale}), together with diaconal service (\latin{diakonia}), so that He Himself continues to teach, sanctify, and shepherd His flock from the pulpit, the confessional, and the altar. Through Matrimony, Christian husband and wife consecrate their home as the ``domestic church'' (\latin{ecclesia domestica}, \emph{Lumen Gentium}~11), imaging in spousal flesh Christ's unbreakable fidelity to His Bride (Eph~5:25--32) and becoming the first heralds of the Gospel to their children by word and example. The public sanctuary and the Christian hearth sustain one another: the altar sends grace into the home, and the home offers sons and daughters to the altar.

\subsection{Christian Marriage and the Counsel of Consecrated Virginity}
\label{subsec:marriage-and-virginity}

Precisely because the Church confesses that marriage is an institution of God Himself and a true sacrament of the New Law, she has always defended its dignity against every dualist or Manichaean contempt for the body. Marriage is ``honourable in all, and the bed undefiled'' (Heb~13:4); procreation belonged to the Creator's blessing in Paradise before sin ever darkened the earth (Gen~1:28; Augustine, \emph{De civitate Dei}~XIV.21). Yet the Church simultaneously guards the evangelical counsel commended by the Lord Himself (Matt~19:12) and preached by St.~Paul (1~Cor~7:32--38): perfect chastity in consecrated virginity or celibacy chosen for the sake of the kingdom of heaven.

Against sixteenth-century rejections of monastic vows and clerical celibacy, the Council of Trent defined the objective ordering of these two holy states in Session~XXIV, canon~10:
\begin{quotation}
\latin{Si quis dixerit, statum coniugalem anteponendum esse statui virginitatis vel caelibatus, et non esse melius ac beatius manere in virginitate aut caelibatu, quam iungi matrimonio: anathema sit.} (If anyone says that the married state is to be preferred to the state of virginity or celibacy, and that it is not better and more blessed to remain in virginity or celibacy than to be joined in matrimony, let him be anathema.)
\end{quotation}
This dogmatic judgment casts no shadow over the goodness of Christian marriage; rather, it illuminates the eschatological meaning of both states (Paul~VI, \emph{Sacerdotalis caelibatus}~19--34; John~Paul~II, \emph{Familiaris Consortio}~16). In the time of earthly pilgrimage, Christian spouses embody the fruitful, bodily union of Christ and His Church through generation and mutual self-gift. Consecrated virgins and celibate ministers anticipate directly the age to come, when earthly marrying ceases and the whole redeemed Church enters forever into the Marriage Supper of the Lamb (Matt~22:30; Apoc~19:7--9). Each vocation illuminates the other: without the holiness of marriage, virginity would lose the spousal sacrifice it transcends; without the sign of virginity, marriage could forget that its deepest horizon lies beyond death in the kingdom of God.
